\chapter{What the Shared Encoder Reads}
\label{ch:encoder-inputs}

One encoder serves every agent, and the same encoder is what a transfer study
moves between grids. Its inputs are raw physical quantities, so before either
form of sharing can be interpreted the inputs themselves have to be measured.
This chapter does that: it establishes what the encoder consumes, how far those
distributions differ between two grids, which of the differences a network
absorbs and which it does not, and what corrections follow.
\\
\\
Chapter~\ref{ch:wcci-transfer} then applies the corrections and runs the
transfer. The measurements here are the reason that study is interpretable at
all: without them, a deficit under transfer could not be told apart from a
representation reading miscalibrated inputs.

\section{Why the input distribution matters}
\label{sec:wcci-why-distribution}

The graph encoder is grid-size independent and can therefore be transferred
across grids. A conventional action head cannot: it has one output per action
index, so its shape changes with the agent partitions and action sets. The
candidate scorer of Section~\ref{sec:candidate-action-scoring} escapes that by
applying shared weights and action-conditioned pooling to a variable set of
actions, which makes the action head transferable too. Only the flat centralized
critic stays grid-specific, its input dimension being set by the grid.
\\
\\
The two input paths must be distinguished. The critic receives the normalized
flat observation; the transferred actor encoder receives only graph node and
edge features, for which physical scaling and running standardization are
disabled here. A frozen \texttt{bus14} encoder therefore reads WCCI's raw graph
quantities, and whether that matters depends on how closely the two
distributions match.


\section{Measurement protocol}
\label{sec:wcci-protocol}

Both environments were rolled out under a do-nothing policy for 6{,}000
environment steps, with a cap of 300 steps per chronic so the sample spreads
over roughly twenty scenarios per grid rather than concentrating in one. At each
step the per-agent local graphs were rebuilt and every visible node recorded,
pooling agents together because a single encoder is shared across them. This
yields 426{,}000 node observations for \texttt{bus14} and 1{,}236{,}000 for
WCCI.
\\
\\
Only nodes are measured, because on this schema only nodes carry measurements:
the nine edge channels are relation indicators, exactly one of which is set on
any active edge, with no continuous content to compare. The transfer runs also
set \texttt{gnn\_include\_neighbors = false}, so every node in a local graph is
controlled and \texttt{domain\_mask} is constant at one throughout.
\\
\\
Four statistics summarize each feature. Writing $F_A$ and $F_B$ for the two
empirical distributions with means $\mu_A,\mu_B$ and standard deviations
$\sigma_A,\sigma_B$, and $\sigma_p=\sqrt{(\sigma_A^2+\sigma_B^2)/2}$ for their
pooled spread:
\begin{align}
  \text{offset} &= \frac{\mu_B-\mu_A}{\sigma_p},
  &\text{scale ratio} &= \frac{\sigma_B}{\sigma_A}, \label{eq:offset-scale} \\
  \mathrm{KS} &= \sup_x |F_A(x)-F_B(x)|,
  &W_1/\sigma_p &= \frac{1}{\sigma_p}\int_0^1 |F_A^{-1}(q)-F_B^{-1}(q)|\,\mathrm{d}q .
\end{align}
{\color{clarif}
\paragraph{Reading the equations.} $A$ is always \texttt{bus14}, the source grid,
and $B$ always WCCI, the target: an offset is therefore positive when WCCI sits
higher, and a scale ratio above one means WCCI is the wider of the two. $F_A$
and $F_B$ are the two \emph{empirical cumulative distribution functions}, one
per grid, built from the samples of Section~\ref{sec:wcci-protocol}: $F_A(x)$ is
the fraction of \texttt{bus14}'s sampled values that fall at or below $x$, so it
rises from $0$ to $1$ as $x$ sweeps the axis. Its inverse $F_A^{-1}(q)$ runs the
other way, returning the value below which a fraction $q$ of the sample lies ---
$F_A^{-1}(0.5)$ is the median. The two definitions are what let the last two
statistics compare the same pair of samples from perpendicular directions: KS
measures the gap between the curves vertically at fixed $x$, $W_1$ between their
inverses horizontally at fixed $q$. Their names: $\mathrm{KS}$ is the
Kolmogorov--Smirnov statistic, the largest vertical gap between the two
empirical cumulative distribution functions. It lies in $[0,1]$, reaching $0$
when the distributions coincide and $1$ when their supports are disjoint, and
the last panel of Figure~\ref{fig:wcci-distributions} draws the gap it measures.
$W_1$ is the
$1$-Wasserstein distance, also called the earth mover's distance: the average
horizontal gap between the two quantile functions, or equivalently the least
mass$\,\times\,$distance needed to rearrange one distribution into the other.
The subscript is the exponent applied to that displacement --- $W_2$ would
square it and so weigh far excursions more heavily --- and this work uses
$p=1$ throughout. Because $W_1$ carries the units of the feature, it is reported
in pooled standard deviations so that channels measured in MW and in degrees
stay comparable.}
\\
\\
The pairing is deliberate: the offset is blind to a change in spread and the
scale ratio is blind to a displacement, so a channel is only unproblematic when
both sit near their neutral values. KS measures disagreement in the bulk and is
comparatively insensitive to tails, which $W_1$ does weigh.
Section~\ref{sec:wcci-interpretation} shows why the two failures need separating.


\section{What the encoder reads}
\label{sec:wcci-features}

The measurements are taken on the disaggregated line-node graph, the schema the
transfer study uses. Its twelve node channels are specified in
Table~\ref{tab:heterogeneous-line-node-features}: four type indicators, the
shared \texttt{p} and \texttt{theta} columns, four line-state columns, the
substation cooldown, and \texttt{domain\_mask}.
\\
\\
Every node type shares one feature vector, so a channel is filled only on the
type that owns it and is zero-filled elsewhere. That makes the sparsity
structural and severe: \texttt{p} and \texttt{theta} carry a value only on
generator and load nodes, which are $24.0\,\%$ of rows on \texttt{bus14} and
$28.7\,\%$ on WCCI, while \texttt{rho} and the other line-state channels are
filled only on transmission-line nodes, $36.6\,\%$ and $36.4\,\%$ of rows. A
statistic taken over all rows therefore mixes a measurement with the zeros of
three other node types, which is why every channel below is reported twice.
\\
\\
Two channels are constant in this sample rather than sparse. The two cooldown
columns are exactly zero only because the policy is do-nothing; under control a
topology change sets the substation cooldown to 3 and a reconnection sets the
line cooldown to 3. \texttt{timestep\_overflow} is zero on all but a handful of
rows. One channel is dimensionless by construction: \texttt{rho} divides each
line's flow by that line's own thermal limit, while every other physical channel
carries the units of the grid it came from.


\section{Results}
\label{sec:wcci-results}

Table~\ref{tab:wcci-shift} reports the shift statistics twice: over all sampled
rows, and over only the node type that carries the channel. The two readings
differ because a shared column is a point mass at zero, contributed by the three
types that do not own it, mixed with the measurement of the type that does.

\begin{table}[H]
  \centering
  \small
  \caption{Distribution shift between the two grids on the disaggregated
  line-node graph, ordered by KS. Offset is in pooled standard deviations and
  the scale ratio is WCCI over \texttt{bus14}. The right-hand block restricts
  each measured channel to the node type that owns it: generator and load nodes
  for \texttt{p} and \texttt{theta}, transmission-line nodes for the line state.
  {\color{clarif}The type indicators and \texttt{domain\_mask} own no type, so
  the restriction does not apply; \emph{constant} marks a channel with no
  variation in this sample, whose ratio is undefined.}}
  \label{tab:wcci-shift}
  \setlength{\tabcolsep}{4pt}
  \begin{tabular}{lrrrrrrrr}
    \toprule
    & \multicolumn{2}{c}{Zeros (\%)} & \multicolumn{3}{c}{All sampled rows}
    & \multicolumn{3}{c}{Owning node type} \\
    \cmidrule(lr){2-3}\cmidrule(lr){4-6}\cmidrule(lr){7-9}
    Channel & \texttt{bus14} & WCCI & offset & ratio & KS & offset & ratio & KS \\
    \midrule
    \texttt{theta}  & 77.5 & 73.8 & $+0.653$ & 0.83 & 0.157 & $+1.864$ & 1.43 & \textbf{0.678} \\
    \texttt{p}      & 77.8 & 76.7 & $-0.037$ & 2.14 & 0.042 & $-0.158$ & \textbf{2.54} & 0.201 \\
    \texttt{rho}    & 63.4 & 63.6 & $-0.088$ & 0.92 & 0.053 & $-0.260$ & 1.05 & 0.141 \\
    \texttt{line\_status} & 63.4 & 63.6 & $-0.005$ & 1.00 & 0.002 & \multicolumn{3}{c}{constant at one} \\
    \texttt{timestep\_overflow} & 100 & 100 & $-0.002$ & 0.70 & 0.000 & $-0.004$ & 0.70 & 0.000 \\
    \texttt{time\_before\_cooldown\_line} & 100 & 100 & --- & --- & 0.000 & \multicolumn{3}{c}{constant} \\
    \texttt{time\_before\_cooldown\_sub}  & 100 & 100 & --- & --- & 0.000 & \multicolumn{3}{c}{constant} \\
    \midrule
    \texttt{is\_busbar}   & 60.6 & 65.0 & $-0.093$ & 0.98 & 0.045 & \multicolumn{3}{c}{---} \\
    \texttt{is\_generator} & 91.5 & 89.3 & $+0.076$ & 1.11 & 0.022 & \multicolumn{3}{c}{---} \\
    \texttt{is\_load}      & 84.5 & 82.0 & $+0.066$ & 1.06 & 0.025 & \multicolumn{3}{c}{---} \\
    \texttt{is\_transmission\_line} & 63.4 & 63.6 & $-0.004$ & 1.00 & 0.002 & \multicolumn{3}{c}{---} \\
    \texttt{domain\_mask}  & 0 & 0 & --- & --- & 0.000 & \multicolumn{3}{c}{constant at one} \\
    \bottomrule
  \end{tabular}
\end{table}

\noindent Read on all rows the table looks harmless: no channel reaches
$\mathrm{KS}=0.2$. That is the type mixing, not agreement. Restricted to the
owning type, \texttt{theta} reaches $0.678$ with $+1.86$ standard deviations of
displacement, and \texttt{p} widens by a factor of $2.54$. The four type
indicators differ only slightly, which says the two grids have a similar mix of
busbars, assets, and lines per local graph rather than saying anything about
their physics.
\\
\\
The extremes make the same point concretely. \texttt{p} spans $[0, 97]$~MW on
\texttt{bus14} and $[-281, 399]$~MW on WCCI, where $7.8\,\%$ of asset rows are
negative and \texttt{bus14} has none. \texttt{theta} spans $[-11.6, 1.5]$
degrees against $[-7.1, 16.8]$: not a wider version of the same distribution but
a displaced one.

\begin{figure}[H]
  \centering
  \includegraphics[width=\linewidth]{figures/wcci_feature_distributions.png}
  \caption{The measured channels on the node type that owns them.
  \texttt{rho} lands almost on top of itself, as expected from a quantity
  already normalized by each grid's own thermal limits. \texttt{p} shares a
  centre but not a width, and WCCI alone produces negative values. \texttt{theta}
  is displaced rather than rescaled: \texttt{bus14} sits almost entirely on the
  negative side. The last panel reads that displacement as the KS statistic.
  Full ranges on log density are in Figure~\ref{fig:wcci-tails-appendix}.}
  \label{fig:wcci-distributions}
\end{figure}


\section{Two failure modes}
\label{sec:wcci-interpretation}

The left panel of Figure~\ref{fig:wcci-offset-scale} separates the channels
along the two axes; the right panel repeats the diagnostic after the corrections
and is read in Section~\ref{sec:wcci-corrected}. The distinction is not
cosmetic, because the architecture responds to the two axes differently.

{\color{clarif}
\paragraph{Reading the axes.} Each point is one channel, placed by the two
statistics of Equation~\eqref{eq:offset-scale} taken over all sampled entries. The horizontal axis
is the \emph{magnitude} of the offset, $|\mu_B-\mu_A|/\sigma_p$, in pooled
standard deviations: how far the two distributions have moved apart, with the
sign discarded because only the size of a displacement matters here. The
vertical axis is the scale ratio $\sigma_B/\sigma_A$, WCCI over \texttt{bus14}:
how much wider or narrower WCCI's spread is, with $1$ meaning identical spread.
It is drawn on a logarithmic scale so that halving and doubling sit the same
distance from $1$, which a linear axis would misrepresent. A channel needing no
correction sits at the origin of this reading --- near-zero offset and ratio
near $1$ --- so distance from the bottom-left corner is distance from
comparability, and the direction of that distance names which of the two
failures it is.}
\\
\\

\begin{figure}[H]
  \centering
  \includegraphics[width=\linewidth]{figures/wcci_feature_offset_scale.png}
  \caption{Absolute offset against scale ratio, over all sampled entries. The
  shaded band and the dashed line mark a scale ratio within a third of unity;
  the dotted vertical line marks an offset of $0.3$ pooled standard deviations.
  \emph{Left:} as measured. Channels outside those bounds fail in one of two
  distinct ways --- the power channels keep their centre and widen, the angles
  keep their width and move. \emph{Right:} the same diagnostic after the two
  corrections of Section~\ref{sec:wcci-correction}, discussed in
  Section~\ref{sec:wcci-corrected}. Points are placed on the owning node type.
  Every channel now sits inside the scale band and the worst displacement falls
  from $1.86$ to $0.81$, but \texttt{p} acquires an offset it did not have.}
  \label{fig:wcci-offset-scale}
\end{figure}

\noindent The actor's node pre-encoder is a linear map followed by ReLU and LayerNorm.
LayerNorm absorbs much of a uniform change of \emph{scale}: multiplying every
input channel by a constant leaves the normalized activations largely unchanged.
It does not absorb an \emph{offset}. Displacing a channel by a constant moves
the pre-activations to a different region of the ReLU, changing which units fire
at all, and normalization applied afterwards cannot undo that.
\\
\\
This makes the angle, not the power, the binding problem. \texttt{p} differs in
width by a factor of $2.54$, which the architecture partially self-corrects,
though the negative WCCI region reaching $-281$~MW is genuinely unseen territory
rather than a wider version of something familiar. \texttt{theta} differs by
$1.86$ pooled standard deviations of displacement, which it does not correct at
all.
\\
\\
\texttt{rho}, meanwhile, needs no correction at all: scale ratio $1.050$, offset
$-0.26$ --- the same values the busbar schema gives, since it is the same set of
lines carrying the same loadings. The single most decision-relevant channel is
already comparable across grids because it was defined in per-unit terms, an
argument for expressing features in grid-relative units wherever the physics
permits it.


\section{Correcting the inputs}
\label{sec:wcci-correction}

Two preprocessing mechanisms exist in the codebase. Deterministic physical
scaling divides each channel by a per-environment physical constant. Running
standardization maintains per-environment, per-node-type statistics over
\texttt{p} and \texttt{theta}, excluding \texttt{rho}, which the measurements
above confirm needs no help. On the diagnosis alone the second looks like the
match, being the only one that can move a displaced channel. It is not.

\subsection{Running standardization makes the mismatch worse}
\label{sec:wcci-running-norm}

Standardizing each environment with its own statistics aligns the aggregate
moments exactly: every channel comes out at offset $0.00$ and scale ratio
$1.00$. It also moves both channels it touches an order of magnitude further
apart (Table~\ref{tab:wcci-running-norm}).

\begin{table}[H]
  \centering
  \small
  \caption{KS between the two grids before and after per-environment
  standardization, over all sampled rows. {\color{clarif}The last column is the
  number a node that does \emph{not} own the channel presents to the encoder
  once standardized. Raw, such a row reads $0$ on both grids; standardizing
  sends it to $-\mu/\sigma$, a property of the grid.} \texttt{rho} is shown for
  contrast and is excluded from the running-statistics feature set.}
  \label{tab:wcci-running-norm}
  \setlength{\tabcolsep}{8pt}
  \begin{tabular}{lrrc}
    \toprule
    Channel & KS raw & KS standardized & ``Not this type'' encoding \\
    & & & \texttt{bus14} vs WCCI \\
    \midrule
    \texttt{p}     & 0.042 & \textbf{0.773} & $-0.371$ vs $-0.145$ \\
    \texttt{theta} & 0.157 & \textbf{0.673} & $+0.481$ vs $-0.143$ \\
    \midrule
    \texttt{rho} (excluded) & 0.053 & 0.634 & $-0.658$ vs $-0.622$ \\
    \bottomrule
  \end{tabular}
\end{table}

\noindent The cause is the shared-column structure of
Section~\ref{sec:wcci-features}. Each channel is a point mass at zero,
contributed by the three node types that do not own it, plus the measurement of
the type that does. Subtracting a per-grid mean moves that point mass to a
grid-dependent location. Since three quarters of all rows are zeros for
\texttt{p} and \texttt{theta}, the encoder's single most frequent input stops
meaning the same thing on the two grids, which is a larger mismatch than the
displacement being corrected. On \texttt{theta} it is starker still: a row that
does not own the channel reads $+0.481$ on \texttt{bus14} and $-0.143$ on WCCI,
so the same absence lands on opposite sides of zero.
\\
\\
Two further properties rule the mechanism out independently of these numbers.
Standardization is scale-invariant, so it erases any scaling applied before it
--- dividing by a constant and then standardizing reproduces standardizing alone
to four decimals --- meaning the two mechanisms cannot be combined. And running
statistics begin empty and converge during training, so a transferred encoder
would spend early target training reading a distribution that is itself still
moving, where a deterministic scaling is fixed at construction and identical on
both grids from the first step.

\subsection{Physical scaling for the power channels}
\label{sec:wcci-physical-scaling}

Physical scaling, defined in Section~\ref{subsec:graph-feature-normalization},
divides the power channels by each environment's own maximum generator capacity.
Two properties answer the failure diagnosed above: being multiplicative it
leaves zeros at zero, so the point mass stays aligned on both grids, and being
fixed at construction it does not drift the way estimated moments do. The
divisor is a single generator's nameplate rating, a crude proxy for grid scale,
so it is worth checking against the alternatives
(Table~\ref{tab:wcci-normalizer}); it is the closest of the four to the divisor
that would equalize the two spreads exactly.

\begin{table}[H]
  \centering
  \small
  \caption{Candidate power normalizers and the \texttt{p} scale ratio each would
  produce on asset nodes. A ratio of $1.00$ would be exact; the raw ratio is
  $2.54$.}
  \label{tab:wcci-normalizer}
  \setlength{\tabcolsep}{8pt}
  \begin{tabular}{lrrr}
    \toprule
    Divisor & \texttt{bus14} & WCCI & Resulting ratio \\
    \midrule
    Exact & --- & --- & 1.00 \\
    \textbf{Maximum $P_{\max}$ (used)} & 140 & 400 & \textbf{0.89} \\
    Total $P_{\max}$ & 540 & 2{,}450 & 0.56 \\
    Generator count & 6 & 22 & 0.69 \\
    Median $P_{\max}$ & 85 & 71 & 3.04 \\
    \bottomrule
  \end{tabular}
\end{table}

\subsection{The angle displacement is a gauge artifact}
\label{sec:wcci-angle-representation}

Scaling cannot repair the angle displacement, and the implementation does not
attempt it: angle channels are divided by a constant $180$, identical on both
grids, so both sides move together and the displacement is untouched. The
measurement explains why. An absolute voltage angle is referenced to whichever
bus is slack, and the two grids do not place theirs alike, so the displacement
in Figure~\ref{fig:wcci-distributions} --- \texttt{bus14} spanning $[-11.6,1.5]$
degrees against WCCI's $[-7.1,16.8]$ --- is largely gauge rather than physics. Nothing that rescales or recentres a reference-dependent quantity
can fix that; the quantity has to be replaced.
\\
\\
The replacement is the angle drop along each line, defined in
Section~\ref{sec:line-node-angle}. Being a difference it is gauge-invariant, and
on this schema it lives on the transmission-line node, which is the most
populous type in the graph, so it also trades a channel owned by a quarter of
the rows for one owned by a third. This study therefore enables it, alongside
physical scaling, in both source and target environments.

\subsection{The action descriptor must not be normalized per agent}
\label{sec:wcci-action-descriptor}

The same error appears once more, on an input this chapter has not otherwise
considered, and with the agents rather than the grids as the populations. One
input reaches the actor without passing through the encoder: the
candidate-action head of Section~\ref{sec:candidate-action-scoring} appends a
static descriptor of each action to the pooled node context. Its count columns
were divided by the largest value in that agent's \emph{own} action set --- the
construction Section~\ref{sec:wcci-running-norm} rejects for node features,
applied to a different input, with the agents as the populations.

\begin{table}[H]
  \centering
  \small
  \caption{Divisors applied to the action count descriptors under the previous
  per-agent rule, \texttt{bus14}.}
  \label{tab:action-descriptor-divisors}
  \setlength{\tabcolsep}{7pt}
  \begin{tabular}{lrrr}
    \toprule
    Agent & Actions & Topology changes & Line endpoints \\
    \midrule
    0 & 57 & 5 & 3 \\
    1 & 51 & 5 & 4 \\
    2 & 83 & 6 & 3 \\
    \bottomrule
  \end{tabular}
\end{table}

\noindent An action moving three elements therefore reached the shared scorer as
$0.60$ for agent~$0$ and $0.50$ for agent~$2$, and every agent's busiest action
reached it as exactly $1.0$ whether it moved five elements or six. On another
grid the divisors change again. The columns are now divided by a constant,
selected by \texttt{candidate\_action\_feature\_scaling}; the busiest actions
then read $1.25$, $1.25$ and $1.5$. This is the only correction here that
applies within a grid as well as across grids, and it qualifies Screen~E of
Chapter~\ref{ch:experiments}, whose runs predate it: their \texttt{f1} cells
appended a descriptor inconsistent between the agents sharing one scorer.

\subsection{The corrected distribution}
\label{sec:wcci-corrected}

Table~\ref{tab:wcci-corrected} measures both changes together, under the same
protocol as Section~\ref{sec:wcci-protocol}, and the right panel of
Figure~\ref{fig:wcci-offset-scale} places the result back on the offset--scale
diagnostic. The angle channel lands in the band occupied by \texttt{rho}, and
the power spread comes within $11\,\%$ of parity.

\begin{table}[H]
  \centering
  \small
  \caption{Distribution shift with \texttt{gnn\_angle\_representation =
  edge\_diff} and \texttt{gnn\_physical\_scaling = true}, against the
  uncorrected baseline, on the node type owning each channel. The offset column
  is included because scaling moves it: dividing each grid by its own constant
  is not offset-preserving when the two constants stand in a different relation
  to typical operation.}
  \label{tab:wcci-corrected}
  \setlength{\tabcolsep}{6pt}
  \begin{tabular}{lrrr|rrr}
    \toprule
    & \multicolumn{3}{c|}{Uncorrected} & \multicolumn{3}{c}{Corrected} \\
    \cmidrule(lr){2-4}\cmidrule(lr){5-7}
    Channel & offset & ratio & KS & offset & ratio & KS \\
    \midrule
    \texttt{theta}      & $+1.864$ & 1.43 & 0.678 & \multicolumn{3}{c}{replaced} \\
    \texttt{theta\_diff} & \multicolumn{3}{c|}{---} & $-0.383$ & \textbf{1.15} & \textbf{0.230} \\
    \texttt{p}          & $-0.158$ & 2.54 & 0.201 & $-0.808$ & \textbf{0.89} & 0.511 \\
    \texttt{rho}        & $-0.260$ & 1.05 & 0.141 & $-0.260$ & 1.05 & 0.141 \\
    \bottomrule
  \end{tabular}
\end{table}

\noindent The angle correction is unambiguous: replacing a reference-dependent quantity
by a gauge-invariant one takes the worst channel in the chapter from $0.678$ to
$0.230$, and its scale ratio stays near one throughout.
\\
\\
The power correction is a genuine exchange rather than a repair. Dividing by
each grid's largest generator rating brings the spread from $2.54$ to $0.89$,
but the raw means, nearly equal at $-0.158$, are not preserved:
\texttt{bus14} runs its generators at $39\,\%$ of its fleet maximum against
WCCI's $20\,\%$, so a divisor taken from the largest unit lands the two grids at
different normalized levels and the displacement grows to $-0.808$. KS rises
with it, from $0.201$ to $0.511$. Section~\ref{sec:wcci-interpretation} argued
that offsets, not scale errors, are what the architecture cannot absorb, so this
moves \texttt{p} onto the axis that binds.
\\
\\
The exchange is still favourable on the evidence available: the worst
displacement across all channels falls from $1.864$ to $0.808$, every scale
ratio ends inside the band, and no channel remains outside both bounds at once.
But it is an exchange, not an elimination. A divisor reflecting typical rather
than peak generation would plausibly do better on both axes;
Table~\ref{tab:wcci-normalizer} tested four candidates against the scale ratio
alone, and selecting one against the offset as well is left open. The transfer
results of Chapter~\ref{ch:wcci-transfer} are the test of whether the residual
displacement matters in practice, and the design of
Section~\ref{sec:wcci-design} runs both input conditions precisely so that
question can be answered.

\begin{figure}[H]
  \centering
  \includegraphics[width=\linewidth]{figures/wcci_feature_corrected.png}
  \caption{Corrected distributions, for comparison with
  Figure~\ref{fig:wcci-distributions}. The angle drop is centred on both grids,
  and the scaled power channels overlap over most of their range instead of
  differing by a factor of $2.5$ in width.}
  \label{fig:wcci-corrected}
\end{figure}


